\section{Reinforcement Learning}
\label{sec:rl}

Reinforcement learning studies how an agent can learn to make sequential decisions through interaction with an environment in order to maximize cumulative reward. At each time step, the agent observes the current state, selects an action, and receives feedback in the form of a reward signal.

Reinforcement learning problems are commonly modeled as a Markov Decision Process (MDP), defined as $M = \langle \States, \Actions, \Transition, \Reward, \Discount \rangle$, where $\States$ is the set of states, $\Actions$ is the set of actions, $\Transition(s' \mid s,a)$ defines the probability of transitioning to state $s'$ after taking action $a$ in state $s$, $\Reward(s,a)$ is the reward function, and $\Discount \in [0,1)$ is the discount factor.

At time step $t$, the agent is in state $s_t \in \States$, selects an action $a_t \in \Actions$, receives reward $r_{t+1} = \Reward(s_t,a_t)$, and transitions to a new state $s_{t+1}$ according to $\Transition(\cdot \mid s_t,a_t)$. The Markov property assumes that the transition and reward depend only on the current state and action.

A policy $\Policy(a \mid s)$ defines the behavior of the agent by specifying a probability distribution over actions given a state. The objective of reinforcement learning is to find a policy that maximizes the expected cumulative reward.

The return from time step $t$ is defined as $G_t = \sum_{k=0}^{\infty} \Discount^k r_{t+k+1}$. The state-value function is defined as $\Value^\Policy(s) = \Expect_\Policy [G_t \mid s_t = s]$, and the action-value function is defined as $\QValue^\Policy(s,a) = \Expect_\Policy [G_t \mid s_t = s, a_t = a]$.

These value functions satisfy recursive relationships known as the Bellman expectation equations. The state-value function satisfies
\[
\Value^\Policy(s) = \Expect_\Policy \left[ r_{t+1} + \Discount \Value^\Policy(s_{t+1}) \mid s_t = s \right],
\]
and the action-value function satisfies
\[
\QValue^\Policy(s,a) = \Expect_\Policy \left[ r_{t+1} + \Discount \Expect_{a' \sim \Policy(\cdot \mid s_{t+1})} \QValue^\Policy(s_{t+1}, a') \mid s_t = s, a_t = a \right].
\]

In discrete settings, the state and action spaces are finite, i.e., $|\States| < \infty$ and $|\Actions| < \infty$. In this case, value functions such as $\Value^\Policy(s)$ and $\QValue^\Policy(s,a)$ can be represented explicitly in tabular form, and reinforcement learning algorithms can directly estimate these quantities for each state or state-action pair.

In continuous settings, the state space $\States \subseteq \Real^n$ and/or the action space $\Actions \subseteq \Real^m$ may be continuous. In such cases, value functions and policies cannot be represented exactly and are instead approximated using parameterized function approximators. A policy may be represented as $\Policy_{\theta}(a \mid s)$, where $\theta$ denotes the parameters of the policy, and a value function may be represented as $\Value_{\phi}(s)$ or $\QValue_{\phi}(s,a)$, where $\phi$ denotes the parameters of the function approximator.

Learning in continuous settings involves optimizing the parameters $\theta$ and $\phi$ to maximize expected return. These formulations are particularly important in robotics and control tasks, where states and actions are naturally continuous and policies must operate over high-dimensional spaces~\cite{goel2022rapidlearn}.


% In every case, we perform a sequence of interactions with the environment, collecting data in the form of trajectories $\tau = (s_0, a_0, r_1, s_1, a_1, r_2, \ldots)$, which are used to update the policy and value function estimates. The goal is to find an optimal policy $\Policy^*$ that maximizes the expected return from any initial state.